\documentclass[11pt]{article}
\usepackage[T1]{fontenc}
\usepackage{lmodern}
\usepackage[margin=1in]{geometry}
\usepackage{amsmath,amssymb,amsthm,booktabs,microtype}
\usepackage[hidelinks]{hyperref}
\newtheorem{theorem}{Theorem}
\newtheorem{lemma}[theorem]{Lemma}
\newcommand{\oeis}[1]{\href{https://oeis.org/#1}{\textup{#1}}}
\title{Sharp position discrepancies for three binary substitutions}
\author{Blake Stone mathematics repository\\[-2pt]\small Proof developed with OpenAI Codex}
\date{November 26, 2026}
\begin{document}
\maketitle
\begin{abstract}
For three of Kimberling's binary fixed points we calculate the infimum and
supremum of the error in the linear approximation to the position of the
$n$-th zero or one. The six pairs of error bounds settle five recorded
inequalities and Dekking's corrected inequality for A284895. In that case the
error $a(n)-n\sqrt2$ lies strictly between 0 and 2 and approaches each
endpoint arbitrarily closely. The proof uses finite equations for extrema of
prefix weights; all computations take place in real quadratic fields.
\end{abstract}

\section{Results}
Each substitution below fixes an infinite word obtained by iterating it on
the stated first letter. Let $p_e(n)$ be the position of its $n$-th letter
$e\in\{0,1\}$, counting both $n$ and positions from one. Assign weight $h$
to zero and $-1$ to one, and define
\[
 D_0(n)=(1+h)n-p_0(n),\qquad D_1(n)=(1+1/h)n-p_1(n).
\]
These use Kimberling's sign convention except that the corrected question
for A284895 is phrased in terms of $-D_1$. In each row of the tables $s$
is the square root indicated for that word.
\begin{center}
\begin{tabular}{@{}lccccc@{}}
\toprule
Word & Seed & $\sigma(0)$ & $\sigma(1)$ & $h$ & $\mu$\\
\midrule
\oeis{A284893}, $s=\sqrt2$ & 0 & 01 & 0111 & $1+s$ & $2-s$\\
\oeis{A284901}, $s=\sqrt3$ & 0 & 01 & 1000 & $s/3$ & $1-s$\\
\oeis{A285668}, $s=\sqrt{17}$ & 1 & 11 & 1101 & $(3+s)/2$ & $(3-s)/2$\\
\bottomrule
\end{tabular}
\end{center}

\begin{theorem}\label{thm:main}
The exact position-error endpoints for all six sequences are as follows.
The last column gives the strict bounds posed in the corresponding entry,
using the sign convention above.
\begin{center}
\begin{tabular}{@{}llccc@{}}
\toprule
Sequence & Error & Infimum & Supremum & Proposed interval\\
\midrule
\oeis{A284894} & $D_0$ & $1+s$ & $3+2s$ & $(2,6)$\\
\oeis{A284895} & $D_1$ & $-2$ & $0$ & $(-2,0)$, corrected\\[2pt]
\oeis{A284902} & $D_0$ & $(-4-2s)/3$ & $(5+2s)/3$ & $(-3,3)$\\
\oeis{A284903} & $D_1$ & $(-6-2s)/3$ & $(9+4s)/3$ & $(-4,6)$\\[2pt]
\oeis{A285669} & $D_0$ & $(-7+s)/16$ & $(29+5s)/16$ & $(-1,4)$\\
\oeis{A285670} & $D_1$ & $(-23+s)/32$ & $(13+5s)/32$ & $(-1,4)$\\
\bottomrule
\end{tabular}
\end{center}
Every indicated strict inequality holds for every $n\ge1$.
\end{theorem}

Dekking proved boundedness for A284894 and A285669, and identified the
slope in the latter case \cite{oeis}. He also observed that A284895 has
$p_1(2)=4$, refuting the original proposed upper bound of one. The
replacement upper bound of two in his comment is optimal. Our argument
uses the classical prefix decomposition of a substitution, as in
Adamczewski \cite{adam}. Jung used a contraction for exact position
endpoints of the different substitution $0\mapsto11,1\mapsto1110$ in
\oeis{A045671} and \oeis{A045672} \cite{jung}.

\section{From positions to finitely many extrema}
Write $w(v)=h|v|_0-|v|_1$ for the weight of a finite word $v$. Direct
substitution in the table of substitutions gives
\begin{equation}
 w(\sigma(v))=\mu w(v),\qquad |\mu|<1.
 \label{eq:scale}
\end{equation}
For example, for A285668 one has $-2=\mu h$ and $h-3=-\mu$. For a letter
$e$, let $B_e$ be the set of weights of prefixes ending \emph{immediately
before} an occurrence of $e$ in the infinite word. For each factorization
$\sigma(a)=v e u$ draw an incoming edge to $e$ with data $(a,w(v))$.
Then
\begin{equation}
 B_e=\bigcup_{(a,d)\to e}(d+\mu B_a).
 \label{eq:sets}
\end{equation}
Indeed every position lies in a substituted block of the fixed word, and
every listed letter in a block contributes a position at each occurrence of
its parent.

Both sets are nonempty. They are bounded because each occurrence lies in
some finite $\sigma^k(\text{seed})$, and tracing it through its ancestors
expresses the prefix weight as a sum of at most $k$ edge offsets times
$1,\mu,\mu^2,\ldots$. If $C$ is the largest absolute offset, this sum has
absolute value at most $C/(1-|\mu|)$.

Put $L_e=\inf B_e$ and $U_e=\sup B_e$. For $\mu>0$, the equation for
$L_e$ takes the minimum of $d+\mu L_a$ over incoming edges, and the
equation for $U_e$ takes the maximum of $d+\mu U_a$. For $\mu<0$ use
$U_a$ in the minimum and $L_a$ in the maximum. These equations have at
most one solution: the minimum/maximum map obeys $\|F(x)-F(y)\|_\infty\le|\mu|\,\|x-y\|_\infty$.
Since the genuine infima and suprema satisfy the equations, checking a
fixed vector proves sharpness as well as enclosure.

\section{The endpoint calculation}
In each list below, $(a,d)$ gives the parent and the offset from the start
of its image. Edges are numbered from one within each list.
\begin{center}\small
\begin{tabular}{@{}lll@{}}
\toprule
Word & Edges into zero & Edges into one\\
\midrule
A284893 & $(0,0),(1,0)$ & $(0,h),(1,h),(1,h-1),(1,h-2)$\\
A284901 & $(0,0),(1,-1),(1,h-1),(1,2h-1)$ & $(0,h),(1,0)$\\
A285668 & $(1,-2)$ & $(0,0),(0,-1),(1,0),(1,-1),(1,h-2)$\\
\bottomrule
\end{tabular}
\end{center}
Substitution in \eqref{eq:sets} gives the following four extrema before
the letter in question.
\begin{center}\small
\begin{tabular}{@{}lcccc@{}}
\toprule
Word & $L_0$ & $L_1$ & $U_0$ & $U_1$\\
\midrule
A284893 & $0$ & $1$ & $2+s$ & $3+2s$\\
A284901 & $(-4-3s)/3$ & $(-1-3s)/3$ & $(5+s)/3$ & $(5+2s)/3$\\
A285668 & $(-31-7s)/16$ & $(-15-7s)/16$ & $(5-3s)/16$ & $(29+5s)/16$\\
\bottomrule
\end{tabular}
\end{center}
For the four columns in order, the attaining edge numbers in their
respective lists are $(1,4,2,2)$, $(2,2,1,1)$, and $(1,4,1,5)$.
These certify equality once the remaining candidates satisfy the
corresponding weak inequalities. All comparisons follow by substituting
the displayed values and using $s^2=2,3,17$, respectively. For example, in
the last row the lower and upper endpoints for zero must be
$-2+\mu U_1$ and $-2+\mu L_1$, since its only edge is from one. For the
lower endpoint for one the five candidates are
$\mu U_0$, $-1+\mu U_0$, $\mu U_1$, $-1+\mu U_1$, and $h-2+\mu U_1$.
Here $U_1>U_0$, $\mu<0$, and $h>2$, giving the fourth candidate.
The supplied exact-arithmetic verifier checks every comparison and the
four attaining equalities in each row.

\begin{proof}[Proof of Theorem~\ref{thm:main}]
If the $n$-th zero has a preceding prefix of weight $x$, the prefix
ending at that zero has weight $x+h=(1+h)n-p_0(n)$. At the $n$-th
one, the analogous weight is $x-1=-hD_1(n)$. It follows that the endpoint
pairs are
\[
 [L_0+h,U_0+h],\qquad [(1-U_1)/h,(1-L_1)/h].
\]
Simplifying yields the table in Theorem~\ref{thm:main}. In five rows both
endpoints are strictly inside the proposed interval. In the remaining
row $D_1(n)=n\sqrt2-p_1(n)$ is irrational for $n\ge1$, so it cannot
attain either integer endpoint. This proves the corrected strict bounds.
\end{proof}

\begin{thebibliography}{9}
\bibitem{oeis} C. Kimberling, OEIS \oeis{A284894}, \oeis{A284895},
\oeis{A284902}, \oeis{A284903}, \oeis{A285669}, \oeis{A285670} (2017).
M. Dekking's notes on A284894 and A284895 date to 2018, and his note on
A285669 to 2021.
\bibitem{adam} B. Adamczewski, \emph{Symbolic discrepancy and self-similar
dynamics}, Ann. Inst. Fourier 54 (2004), 2201--2234.
\href{https://doi.org/10.5802/aif.2079}{doi:10.5802/aif.2079}.
\bibitem{jung} J. M. Jung, comments on \oeis{A045671} and \oeis{A045672},
November 18, 2026, including sharp extrema for the $1110$ substitution.
\end{thebibliography}
\end{document}
